% ABSTRACT (English)
% Progress figures updated to match Chapter 5 as of the current report (status as of August 2026).
% The citation header (authors/year/advisor) and the Keywords line are generated
% automatically from main.tex; do not add them in this file.

This project aims to develop the web application "Dino" as a centralized tourism information platform in Khon Kaen province to address the problem of fragmented and incomplete tourism data. The project integrates artificial intelligence (AI) technologies via Large Language Models (LLMs) combined with Retrieval-Augmented Generation (RAG) architecture to develop two main systems: (1) an Adaptive Trip Planning system that generates itineraries based on constraints regarding budget, time, and personal interests, with plan quality automatically verified through an LLM-as-a-Judge mechanism, and (2) a real-time Chatbot system for interactive communication and tourism recommendations. The system utilizes a database of attractions and restaurants retrieved from the Google Places API within Khon Kaen province.

Current results indicate that the backend, frontend, and AI service are fully connected and operating end-to-end, including the streaming chatbot and the LLM-based trip planner, with an estimated 85--90\% of the designed scope complete. However, user authentication and the QR-based points collection/redemption system remain prototype-level features that are not yet connected to persistent storage, and the system has not yet undergone broad User Acceptance Testing (UAT). Once the remaining components are completed and tested with real users, this application is expected to reduce information search time, enhance the efficiency of personalized trip planning, support local entrepreneurs, and promote the sustainable tourism economy of Khon Kaen province.
